% Part: proof-theory
% Chapter: sequent-calculus
% Section: rules-G3c

\documentclass[../../../include/open-logic-section]{subfiles}

\begin{document}

\begin{table}
\begin{tabular}{@{}c@{}c@{}}
\multicolumn{2}{@{}c@{}}{అక్షయాలు}\\[2ex]
$!C, \Gamma \Sequent !C$ & $\lfalse, \Gamma \Sequent \Delta$
\\[2ex]
\multicolumn{2}{@{}c@{}}{తార్కిక నియమాలు}\\[2ex]
\Axiom$\lnot !A, \Gamma \fCenter !A $
\RightLabel{\LeftR{\lnot}}
\UnaryInf$\lnot !A, \Gamma \fCenter \Delta$
\DisplayProof
&
\Axiom$!A, \Gamma \fCenter$
\RightLabel{\RightR{\lnot}}
\UnaryInf$ \Gamma \fCenter \lnot !A $
\DisplayProof
\\[4ex]
\Axiom$ !A, !B, \Gamma \fCenter \Delta$
\RightLabel{\LeftR{\land}}
\UnaryInf$ !A \land !B, \Gamma \fCenter \Delta$
\DisplayProof
&
\Axiom$\Gamma \fCenter !A$
\Axiom$ \Gamma \fCenter !B$
\RightLabel{\RightR{\land}}
\BinaryInf$ \Gamma \fCenter !A \land !B $
\DisplayProof
\\[4ex]\Axiom$!A, \Gamma \fCenter \Delta$
\Axiom$!B, \Gamma \fCenter \Delta$
\RightLabel{\LeftR{\lor}}
\BinaryInf$!A \lor !B, \Gamma \fCenter \Delta$
\DisplayProof
&
\Axiom$\Gamma \fCenter !A_i$
\RightLabel{\RightR{\lor}}
\UnaryInf$ \Gamma \fCenter !A_1 \lor !A_2$
\DisplayProof
\\[4ex]
\Axiom$!A \lif !B, \Gamma \fCenter !A$
\Axiom$!B, \Gamma \fCenter \Delta$
\RightLabel{\LeftR{\lif}}
\BinaryInf$!A \lif !B, \Gamma \fCenter \Delta$
\DisplayProof
&
\Axiom$!A, \Gamma \fCenter !B$
\RightLabel{\RightR{\lif}}
\UnaryInf$ \Gamma \fCenter !A \lif !B $
\DisplayProof
\\[4ex]
\Axiom$ !A(t), \lforall[x][!A(x)], \Gamma \fCenter \Delta$
\RightLabel{\LeftR{\lforall}}
\UnaryInf$ \lforall[x][!A(x)],\Gamma \fCenter \Delta$
\DisplayProof
&
\Axiom$ \Gamma \fCenter !A(a) $
\RightLabel{\RightR{\lforall}}
\UnaryInf$ \Gamma \fCenter \lforall[x][!A(x)]$
\DisplayProof
\\[4ex]
\Axiom$ !A(a), \Gamma \fCenter \Delta $
\RightLabel{\LeftR{\lexists}}
\UnaryInf$ \lexists[x][!A(x)], \Gamma \fCenter \Delta$
\DisplayProof
&
\Axiom$ \Gamma \fCenter !A(t) $
\RightLabel{\RightR{\lexists}}
\UnaryInf$ \Gamma \fCenter \lexists[x][!A(x)]$
\DisplayProof
\end{tabular}
\caption{అంతఃప్రజ్ఞావాద తర్కానికి \Log{G3i} నియమాలు. $\RightR{\forall}$, $\LeftR{\exists}$లలో \tetoken{స్థిరాంకం}{!!{constant}}~$a$ నిర్ధారణ సీక్వెంట్‌లో రాకూడదు. $\Gamma$ \tetoken{సూత్రాల}{!!{formula}s} బహుసమితి; $\Delta$ ఖాళీగా లేదా ఒకే \tetoken{సూత్రంతో}{!!{formula}} ఉండవచ్చు. అక్షయాల్లో \tetoken{సూత్రం}{!!{formula}}~$!C$ పరమాణువు కావాలి. \Log{G1i}లో \RightR{\Weakening}ను, అక్షయం $\lfalse \Sequent \quad$ను తీసేస్తే \Log{G1m} వస్తుంది.}
\ollabel{tab:G3i}
\end{table}
\sourcecorrection{OLTEPTSEQG3I-001}{మూల వియోగ కుడి నియమం ఆధారంలో రెండు ఉత్తరపక్ష సూత్రాలు ఉండటం ఈ అంతఃప్రజ్ఞావాద పట్టిక పరిమితికి విరుద్ధం. $i=1,2$లకు ఒక్కొక్క భాగం ఆధారంగా రెండు రూపాలను సూచించే ఒకే పథకంగా సరిచేశాం.}
\sourcecorrection{OLTEPTSEQG3I-002}{సార్వత్రిక ఎడమ ఆధారంలో ప్రధాన సూత్రం, పక్క సందర్భం మధ్య కామాను చేర్చాం.}
\sourcecorrection{OLTEPTSEQG3I-003}{పక్క గమనికలో మొదటి అంతఃప్రజ్ఞావాద వ్యవస్థ కనిష్ఠ రూపం కోసం మూలం అసత్య అక్షయాల తొలగింపును మాత్రమే చెప్పింది. ముందరి పట్టికలోని పూర్తి నిర్వచనం ప్రకారం కుడి బలహీనీకరణ తొలగింపును కూడా స్పష్టం చేసి, అసలు ఆధార అసత్య అక్షయాన్ని రాశాం.}
\sourcecorrection{OLTEPTSEQG3I-004}{మూల పట్టిక గుర్తు సాంప్రదాయ వ్యవస్థ గుర్తునే మళ్లీ వాడింది; ఈ అంతఃప్రజ్ఞావాద పట్టికకు ప్రత్యేక గుర్తు ఇచ్చాం.}

\end{document}
